% Generated draft: requires mathematical and source review.
\documentclass[11pt]{article}
\usepackage[margin=25mm]{geometry}
\usepackage{fontspec}
\setmainfont{Latin Modern Roman}
\usepackage{amsmath,amssymb,mathtools,unicode-math}
\setmathfont{Latin Modern Math}
\usepackage{enumitem,needspace,xurl,hyperref}
\setlength{\parindent}{0pt}
\setlength{\parskip}{5pt}
\setlength{\emergencystretch}{3em}
\newcommand{\N}{\mathbb{N}}
\newcommand{\Z}{\mathbb{Z}}
\newcommand{\Q}{\mathbb{Q}}
\newcommand{\R}{\mathbb{R}}
\newcommand{\C}{\mathbb{C}}
\newcommand{\F}{\mathbb{F}}
\newcommand{\A}{\mathbb{A}}
\newcommand{\PP}{\mathbb{P}}
\newcommand{\E}{\mathbb{E}}
\newcommand{\eps}{\varepsilon}
\newcommand{\dd}{\,\mathrm{d}}
\begin{document}
\section*{UnsolvedMath ID: 800224}
\textbf{Source problem code:} catalogue entry
\subsection{Definitions and mathematical statement}
A domain in \(\C^n\) is a nonempty connected open set. To specify a peak-set problem, fix a domain \(D\subseteq\C^n\) and an algebra
\[
\mathcal A\subseteq\{f\in C(\overline D): f|_D\in\mathcal O(D)\},
\]
where \(\mathcal O(D)\) denotes the holomorphic functions on \(D\). Under a boundary peak-set convention, a closed set \(E\subseteq\partial D\) is an \(\mathcal A\)-peak set if there exists \(f\in\mathcal A\) such that
\[
f(z)=1\quad(z\in E),\qquad |f(z)|<1\quad(z\in\overline D\setminus E).
\]
For example, one possible choice of algebra is
\[
A(D)=\{f\in C(\overline D):f|_D\in\mathcal O(D)\}.
\]
This example is an explicit possible interpretation, not an algebra specified by the supplied source.

The supplied problem asks: describe peak sets on weakly pseudoconvex domains of finite type. With the preceding boundary convention and an algebra \(\mathcal A\) fixed, the requested characterization would determine precisely which closed sets \(E\subseteq\partial D\) satisfy the displayed peak-function condition as \(D\) ranges over the intended class of domains.

The source does not specify the function algebra, whether peak sets are required to lie in the boundary, the complex dimension, boundedness, boundary regularity, or the definition and locus of finite type. It also does not specify whether ``weakly pseudoconvex'' requires Levi degeneracy somewhere or merely allows it. Thus the displayed definition provides a conditional formulation; the supplied sentence does not determine a unique fully specified classification problem.
\subsection{Short English statement}
Characterize the sets on which an admissible holomorphic function equals \(1\) while having modulus strictly less than \(1\) everywhere else, for weakly pseudoconvex domains of finite type. The admissible function class and the precise domain and finite-type hypotheses remain unspecified in the supplied source.
\subsection{Sources}
\begin{itemize}
\item Local UnsolvedMath registry, permanent ID \(800224\), exact source problem code \texttt{catalogue entry}; title: \emph{Several Complex Variables Problem List --- Problem 20}. The supplied registry labels the problem \texttt{open}; this is a supplied status label, not an independently verified mathematical assessment.
\item Buzzard, Fornaess, and Krantz, \emph{The Several Complex Variables Problem List} (1995), Problem 20. Supplied source URL: \url{https://arxiv.org/abs/math/9509202}. The registry attributes the problem to these authors and records proposed year 1995; no separate individual proposer is identified.
\item Supplied provenance records access on 2026-10-16 and extraction method \texttt{arxiv\_latex}. These are supplied provenance details; the source has not been independently visited or verified for this draft.
\end{itemize}
\end{document}
